\section{Related work}

\paragraph{Networked agent infrastructure and protocols.}
Agent-to-agent task protocols publish a capability card and accept tasks while
deliberately keeping a counterpart opaque \cite{a2a2025}; broadcast networks for agents
route published needs and capabilities among strangers while stating that a
participant's private information does not enter the network \cite{eigenflux}; surveys
catalogue the wider protocol design space \cite{agentprotocols2025,coral2025}. These
systems encode precisely the split we intervene on---whether an edge carries a
delegated task or direct access to state---but to our knowledge none has been evaluated
for which side of it changes collective behavior.

\paragraph{Discovery, capability representation, and topology.}
Capability discovery at network scale is the first thrust of current agentic-web
research, and agent cards are its dominant representation. Our formation axis is an
intervention on exactly that: which cards a requester may enumerate and address. We
treat card contents as an experimental variable rather than an interface detail, and
find that the addressable set matters far more than what the cards say.

\paragraph{Trust, reputation, and coordination.}
Reputation systems for agent networks assume relational metadata conditions behavior.
Our attribute axis tests that assumption directly and does not find support for it at
the level where such metadata is usually carried, which bears on whether
manipulation-resistant reputation is solving the binding constraint. We are explicit
that this bounds asserted metadata only, not enforcement-coupled reputation.

\paragraph{Human ties, teams, and information pooling.}
The intuition that open networks supply variety and bounded ones supply follow-through
descends from work on weak ties \cite{granovetter1973}, exploration and exploitation
\cite{march1991}, transactive memory \cite{wegner1987}, and open innovation
\cite{chesbrough2003}; that literature studies humans, where the properties in our
attribute group are real rather than asserted. Our task construction is the
hidden-profile paradigm \cite{stasser1985}, adopted because it supplies exact ground
truth and makes pooling rather than retrieval the object of measurement.

\paragraph{Capabilities and evaluation.}
The authority model is object-capability discipline \cite{miller2006}, so a condition
is a set of grants a kernel checks rather than a convention an agent could ignore.
Where an answer key cannot reach, blind pairwise comparison with a cross-family judge
is the standard remedy \cite{zheng2023}, since evaluators favour their own generations
\cite{panickssery2024}; we preregistered such a protocol but did not run it, and no
result here depends on a judge.
